\documentclass[11pt]{article}
\usepackage{preamble}
\renewcommand{\answer}[1]{}

\begin{document}

\ladderheading{AP Statistics \textperiodcentered\ Unit 1 \textperiodcentered\ Topic 1.13 \textperiodcentered\ B: 2026-10-05 \textperiodcentered\ E: 2026-10-16}{When Blocking Earns Its Keep}

\focusbanner{TODAY'S PROBLEM}

Suppose a grower has 24 similar tomato plants, 12 in a sunny greenhouse and 12 in a shaded greenhouse. The response is fruit yield in kilograms per plant.\par\smallskip \textbf{Design A:} Completely randomize all 24 plants; assign 12 to fertilizer and 12 to no fertilizer, ignoring greenhouse.\par \textbf{Design B:} Treat greenhouse as a block; within each greenhouse, randomly assign 6 plants to fertilizer and 6 to no fertilizer.\par\smallskip Pilot means from plants grown under the same conditions are shown.\par
\begin{center}
\textbf{Mean fruit yield (kg per plant)}\par\smallskip
\begin{tabular}{|l|c|c|}
\hline
\textbf{Greenhouse} & \textbf{Fertilizer} & \textbf{None} \\ \hline
\textbf{Sunny} & 7.2 & 6.0 \\ \hline
\textbf{Shaded} & 4.0 & 2.8 \\ \hline
\end{tabular}
\end{center}
\begin{firsttakebox}{FIRST TAKE --- before the video (5 min)}
Which design would you use to detect the fertilizer effect most clearly? Name one detail behind your choice.
\workspace{0.9in}
\end{firsttakebox}

\begin{callout}{\IconBook\ BLOCK, RANDOMIZE, THEN INFER (keep on the page)}{calloutgreen}
In a \textbf{completely randomized design}, treatments are assigned across all units. In a
\textbf{randomized block design}, first group units that are similar on a relevant variable, then randomly
assign treatments \emph{within each block}. Blocking separates natural block-to-block variability from the
treatment comparison. Random assignment supports cause only when the observed difference is too large for
chance alone; generalization also requires units that represent the target population.
\end{callout}

\newpage
\textbf{Finish after the video:}\par\smallskip

\dokbadge{1}~\textbf{(a)} Identify the treatments, response variable with units, and blocking variable in Design B.\par
\workspace{0.8in}

\dokbadge{2}~\textbf{(b)} Compute the fertilizer-minus-control difference within each greenhouse and the sunny-minus-shaded difference within each treatment.\par
\workspace{1.4in}

\dokbadge{3}~\textbf{(c)} Which design should the grower use? Cite the relative sizes of the greenhouse and fertilizer differences that settle your choice, and explain what chance imbalance could do to Design A. If Design B produced a statistically significant fertilizer difference, state the causal conclusion it would permit and one limit on generalizing it.\par
\workspace{2.0in}

\begin{sentenceframebox}\raggedright\textbf{Frame:} Design \blank[0.4in] is stronger because the block difference is \blank[0.7in] kg versus a treatment difference of \blank[0.7in] kg.\end{sentenceframebox}

\begin{sentenceframebox}\raggedright\textbf{Frame:} If significant, the result supports \blank[2.2in] for \blank[1.6in], but not \blank[1.5in].\end{sentenceframebox}

\begin{summaryexitbox}{EXIT --- turn this in}
One thing the video changed about my first take: \hrulefill\par
\workspace{0.4in}
\end{summaryexitbox}

\end{document}
